\section{An application: factorization uniformly in the degree}
\label{sec:fta-extension}

The repository's one-root FTA proof is uniform in the numerical degree,
but its complete-factorization paragraph deliberately claims only a
fixed-degree schema.  It asks for an additional bound on nested precision
requests before treating a variable number of deflations.  We supply that
bound here.  The input coefficient evaluator remains one fixed primitive
recursive program.  The result does not evaluate arbitrary codes of
primitive recursive programs.

The only substantive algebraic dependency is the repository's finite
rational descent lemma, whose proof uses bounded radius selection, five
fixed-length square-root bisections, and a finite binomial polynomial
\cite{repo-fta}.  We state its complete contract, reproduce the monic
one-root argument, and prove all additional estimates needed for uniform
factorization.  This is a human mathematical extension of the repository
proof, not an assertion that its outstanding deeply embedded Lean
derivation has already been constructed.

\subsection{Norms, finite names, and rounding}

For a Gaussian rational $z=x+iy$, put $|z|_1=|x|+|y|$.
Then $|zw|_1\le |z|_1|w|_1$, $|z|\le |z|_1\le 2|z|$, where
$|z|$ is the Euclidean norm.  Euclidean inequalities in the descent
argument are expressed arithmetically by squared rational comparisons.
For a finite polynomial $F$, put
\[
 \|F\|_1=\sum_j |[X^j]F|_1.
\]
Finite convolution and the triangle inequality prove
$\|FG\|_1\le\|F\|_1\|G\|_1$.

Let $\operatorname{rd}_v(z)$ round each real coordinate down to the
nearest multiple of $2^{-v}$.  This is a total primitive recursive
Gaussian-rational operation, and
\[
 |\operatorname{rd}_v(z)-z|_1\le 2^{1-v}.
 \tag{F1}\label{eq:fta-round}
\]
The convention for negative coordinates is ordinary rational floor.
Its error estimate follows directly from the two division-with-remainder
inequalities.  In particular $\operatorname{rd}_v(1)=1$.

For $n\ge1$ and $M\ge2$, call a coefficient table
$a_i(s)$, $0\le i\le n$, $0\le s\le h$, good through $h$ if
\begin{align*}
 a_n(s)&=1, & |a_i(s)|_1&\le M,\\
 |a_i(s)-a_i(t)|_1&\le 2^{-s}+2^{-t}
 &&(i\le n,\ s,t\le h).
\end{align*}
All quantifiers in this definition are bounded.  Its characteristic
function is a finite Boolean conjunction, so goodness is expressed by
one primitive recursive equation.  An infinite name is merely convenient
notation for the corresponding family of finite tables.

The use of $|\cdot|_1$ and an explicit coefficient cap is harmless for
the repository's fast Euclidean names.  Given such names $a_i$, replace
the lower coefficients by
\[
 \widetilde a_i(s)=\operatorname{rd}_{s+4}(a_i(s+3)),
 \qquad \widetilde a_n(s)=1.
\]
The new names represent the same coefficients and satisfy the displayed
$|\cdot|_1$ regularity.  Indeed the two source-name errors contribute at
most $\tfrac14(2^{-s}+2^{-t})$, and rounding contributes at most
$\tfrac18(2^{-s}+2^{-t})$.  A suitable cap is
\[
 M=4+\left\lceil\max_{i<n}|a_i(0)|_1\right\rceil.
\]
Here monicity is supplied as an exact presentation of the top coefficient;
we do not try to decide whether an arbitrary coefficient name equals $1$.

\subsection{A uniform finite root routine with an explicit query bound}

The following finite rational lemma is the imported result.

\begin{lemma}[Finite rational descent contract]\label{lem:fta-descent-contract}
For $n\ge1$, put $D=2^n4^{2n^2}$ and $q=1-(4D)^{-1}$.
There is one primitive recursive routine
$\operatorname{Descent}(n,P,z)$ such that for every nonconstant Gaussian
rational polynomial $P$ of degree at most $n$, it returns Gaussian rational
$\Delta$ and nonnegative rational $r_0$ with
\[
 |P(z+\Delta)|\le q|P(z)|,\qquad |\Delta|\le r_0,
 \qquad |b_i|r_0^i\le |P(z)|\quad(1\le i\le n),
\]
where $P(z+W)=\sum_i b_iW^i$.
All loops and searches in this routine have explicit primitive recursive
bounds in the input codes.  Failed searches and malformed inputs have
fixed default outputs.
\end{lemma}

For clarity, the bound is uniform in $n$; it is not an external family
of routines indexed by fixed $n$.  The source proves this lemma by
testing a finite range of dyadic radii, finding a repeated maximizing
index among $2n+1$ scales, and constructing the appropriate direction by
the finite binomial identity.  These operations introduce no new calls
to the coefficient-name evaluator.

Fix $n,M$ as above and define
\begin{align*}
 B&=M+2, &D&=2^n4^{2n^2}, &K&=8nD,\\
 \rho&=1-K^{-1}, &Q&=1-(8D)^{-1}, &R&=BK+1,\\
 L&=2+\operatorname{bl}(16D(n+1)R^n),
 &H&=\operatorname{bl}(BK),\\
 J(t)&=K(t+H+2), &U(n,M,t)&=L+J(t).
 \tag{F2}\label{eq:fta-constants}
\end{align*}
Here $\operatorname{bl}(b)$ is the binary digit count, so
$b<2^{\operatorname{bl}(b)}$ for $b>0$.
The symbols $B,D,K,R,L,H,J$ in subsequent formulas refer to these explicit
functions of $n,M$.

For the given coefficient table set
\[
 P_j(X)=\sum_{i=0}^n a_i(L+j)X^i,
 \qquad z_0=0,
 \qquad z_{j+1}=z_j+\operatorname{Descent}(n,P_j,z_j).\Delta.
 \tag{F3}\label{eq:fta-trajectory}
\]
This is one fixed primitive recursion.  Define
\[
 \alpha(t)=z_{J(t)},\qquad
 r(t)=\operatorname{rd}_{t+4}(\alpha(t+2)),\qquad A=2R+1.
 \tag{F4}\label{eq:fta-root-normalized}
\]

\begin{lemma}[Finite monic root routine]\label{lem:fta-root-table}
For an output horizon $h$, goodness of the input table through
$U(n,M,h+2)$ suffices to prove
\begin{align*}
 |r(t)|_1&\le A &&(t\le h),\\
 |r(t)-r(u)|_1&\le 2^{-t}+2^{-u} &&(t,u\le h).
\end{align*}
The value $r(t)$ reads coefficient entries only at indices at most
$U(n,M,t+2)$, and it is unchanged if the input table is extended.
In addition, setting
\[
 E=4(n+1)^2(M+1)(A+1)^{n+1},
 \tag{F5}\label{eq:fta-E}
\]
the same finite hypotheses give
\[
 \left|\sum_{i=0}^n a_i(t)r(t)^i\right|_1\le E2^{-t}
 \qquad(t\le h).
 \tag{F6}\label{eq:fta-same-precision-residual}
\]
Every assertion has a quantifier-free-induction proof after the displayed
finite horizons are fixed.
\end{lemma}

\begin{proof}
The leading coefficient of every $P_j$ is $1$.  The radius certificate
in Lemma~\ref{lem:fta-descent-contract} therefore gives
$r_0^n\le |P_j(z_j)|$.  Finite Bernoulli arithmetic gives
$Q\le\rho^n$, and hence
\[
 |P_j(z_j)|\le BQ^j\quad\Longrightarrow\quad
 |\Delta_j|\le B\rho^j.
\]
Indeed $(B\rho^j)^n\ge B\rho^{nj}\ge BQ^j$, and positive integer
powers reflect order on nonnegative rationals.  This argument avoids
computing any $n$th root.

Induct simultaneously on the finite inequalities
\[
 |P_j(z_j)|\le BQ^j,
 \qquad |z_j|\le B\sum_{v<j}\rho^v\le BK<R.
 \tag{F7}\label{eq:fta-monic-invariant}
\]
The base case uses $|a_0(L)|\le M\le B$.  The step bound just proved
establishes the position bound at $j+1$.  At any point $|x|\le R$, the
adjacent coefficient-name comparisons give
\[
 |P_{j+1}(x)-P_j(x)|
 \le\frac32(n+1)R^n2^{-L-j}
 \le\frac{3}{32D}2^{-j}
 \le\frac{1}{8D}BQ^j.
\]
The choice of $L$ is stronger than the inequality used here.  Rational
descent then changes the residual bound by the factor
$q+(8D)^{-1}=Q$.  This proves \eqref{eq:fta-monic-invariant}.

Finite geometric summation gives
\[
 |z_v-z_u|\le BK\rho^u\qquad(v\ge u).
\]
Also $\rho^K\le\tfrac12$: apply finite Bernoulli to
$(1+(K-1)^{-1})^K$.  Thus
\[
 |\alpha(v)-\alpha(u)|\le 2^{-u-2}\qquad(v\ge u).
 \tag{F8}\label{eq:fta-raw-root-cauchy}
\]
If $s\ge U(n,M,t)$ and the input table also contains row $s$, comparison
with $P_{J(t)}$ gives
\[
 \left|\sum_{i=0}^n a_i(s)\alpha(t)^i\right|
 \le BQ^{J(t)}+2(n+1)R^n2^{-L-J(t)}
 \le 2^{-t}.
 \tag{F9}\label{eq:fta-raw-root-residual}
\]
This is a finite statement about the specified $s$, not a limit argument.

Rounding in \eqref{eq:fta-root-normalized} adds $|\cdot|_1$ error at
most $2^{-t-3}$.  Combining this with twice the Euclidean bound
\eqref{eq:fta-raw-root-cauchy} proves regularity of $r$ (in fact with
slack), and $|r(t)|_1\le 2R+1=A$.

For \eqref{eq:fta-same-precision-residual}, use
\eqref{eq:fta-raw-root-residual} at $t+2$ and
$s=U(n,M,t+2)$.  Replacing row $s$ by row $t$ costs at most
$2(n+1)A^n2^{-t}$ in $|\cdot|_1$.  The residual in
\eqref{eq:fta-raw-root-residual} costs at most $2^{-t-1}$ in that norm.
Finally replacing $\alpha(t+2)$ by $r(t)$ costs at most
\[
 \frac18 M\left(\sum_{i=1}^n iA^{i-1}\right)2^{-t}
\]
by the finite difference-of-powers identity.  Their sum is less than
$E2^{-t}$ for \eqref{eq:fta-E}.

The largest coefficient index used in these arguments is the displayed
$U(n,M,h+2)$.  In the trajectory, the invariant is a pair of rational
inequalities; its Euclidean components can be squared.  Geometric-sum
and coefficient comparisons are finite folds.  A Boolean conjunction
stores all required assertions through the fixed horizon.  Quantifier-free
induction therefore proves the finite contract.  Prefix compatibility
follows because each run follows exactly the recursion
\eqref{eq:fta-trajectory}, with no decision based on table length.
\end{proof}

\subsection{A normalized deflator and its complete finite contract}

Use the following additional integer bounds:
\begin{align*}
 M^+&=2+(n+1)(M+1)(A+1)^n,\\
 a&=\operatorname{bl}(4E),\\
 D^+&=4(n+1)(M^++1)(A+1),\\
 V(n,M,h)&=U(n,M,h+a+2).
 \tag{F10}\label{eq:fta-deflation-budgets}
\end{align*}
The letter $a$ in this display denotes the precision shift, not a
coefficient.  For $0\le j<n$ put
\begin{align*}
 q_j(v)&=\sum_{\ell=j+1}^n a_\ell(v)r(v)^{\ell-j-1},\\
 b_j(t)&=\operatorname{rd}_{t+4}(q_j(t+a)),\\
 P(t;X)&=\sum_{j=0}^n a_j(t)X^j,
 &P^+(t;X)&=\sum_{j=0}^{n-1} b_j(t)X^j.
 \tag{F11}\label{eq:fta-deflator}
\end{align*}
These are finite sums, and $b_{n-1}(t)=1$ exactly.

\begin{lemma}[Finite deflation contract]\label{lem:fta-deflation-table}
If the input coefficient table is good through $V(n,M,h)$, then the
output table $b_j(t)$ is monic, bounded by $M^+$ in $|\cdot|_1$, and
regular through $h$.  The normalized root is regular and bounded by $A$
through $h$.  For every $t\le h$,
\[
 \|P(t;X)-(X-r(t))P^+(t;X)\|_1\le D^+2^{-t}.
 \tag{F12}\label{eq:fta-defect}
\]
The complete calculation and every finite verification use only input
indices at most $V(n,M,h)$.  Extending the input table changes none of the
already computed output entries.
\end{lemma}

\begin{proof}
Each $q_j$ has at most $n$ terms.  Hence
\[
 |q_j(v)|_1\le nM(A+1)^n<M^+-1.
\]
Rounding contributes at most $2^{-t-3}\le1$, proving the cap.
For any two available indices $v,w$, subtract corresponding terms and
apply the difference-of-powers identity.  Input and root regularity give
\[
 |q_j(v)-q_j(w)|_1
 \le\left(\sum_{r<n} A^r+M\sum_{r<n}rA^{r-1}\right)
       (2^{-v}+2^{-w})
 \le E(2^{-v}+2^{-w}),
\]
where the $r=0$ term of the second sum is $0$.
Since $E2^{-a}<1/4$, after the index shift and rounding this is at most
$\tfrac38(2^{-t}+2^{-u})$ for $b_j(t),b_j(u)$, which implies the
required regularity.

Set $v=t+a$ and $Q_v(X)=\sum_{j<n}q_j(v)X^j$.  Synthetic division is
the exact finite identity
\[
 P(v;X)=(X-r(v))Q_v(X)+P(v;r(v)).
\]
It follows by cancelling adjacent coefficients, including the constant
coefficient.  The defect in \eqref{eq:fta-defect} is the sum of four
terms:
\begin{align*}
 &[P(t;X)-P(v;X)] + P(v;r(v))\\
 &\quad +(X-r(v))[Q_v(X)-P^+(t;X)]
       +(r(t)-r(v))P^+(t;X).
\end{align*}
Their coefficient norms are bounded respectively by
\[
 2(n+1)2^{-t},\qquad
 E2^{-v}<\tfrac14 2^{-t},\qquad
 (A+1)n2^{-t-3},\qquad
 2nM^+2^{-t}.
\]
Their sum is at most $D^+2^{-t}$.

The residual and root regularity at indices at most $h+a$ require the
source root contract through $U(n,M,h+a+2)=V(n,M,h)$.  All other
coefficient queries are smaller.  Thus $V$ is a verification bound as
well as an evaluation bound.  The proof consists of finite coefficient
identities, bounded sums, and rational comparisons; no induction on a
statement about all precisions has been used.
\end{proof}

\subsection{Variable-degree iteration by finite tables}

Let the input degree be $d\ge1$ and initial cap be $M_0=M$.
Compute the numeric caps and constants by the single primitive recursion
\[
 n_i=d-i,\qquad
 M_{i+1}=M^+(n_i,M_i)\quad(0\le i<d).
 \tag{F13}\label{eq:fta-cap-recursion}
\]
Write $A_i,D_i^+,V_i$ for the corresponding functions in
\eqref{eq:fta-deflation-budgets}.  No coefficient evaluation is needed
to compute these bounds.

For a desired output horizon $t$, determine input horizons backwards:
\[
 h_d=t,\qquad h_i=V_i(h_{i+1})\quad(i=d-1,\ldots,0).
 \tag{F14}\label{eq:fta-backwards}
\]
This is a primitive recursion over a reversed finite index list.  It is
not recursion on a program code.  Each $V_i$ is monotone in its last
argument and is at least that argument.

Generate the original coefficient table through $h_0$ by the one fixed
PR input evaluator.  At stage $i$, run the finite deflator on its current
table to obtain its root entries and next coefficient entries through
$h_{i+1}$.  After $d$ stages the residual polynomial is exactly $1$.
Each stage is one application of the same fixed PR table routine to
numeric inputs $n_i,M_i,h_{i+1}$ and the current finite table.  Therefore
the entire staged calculation is primitive recursive uniformly in $d$.

For an output root at index $i<d$ and precision $p$, use this algorithm
with terminal horizon $p$ and select the $i$th root entry at $p$.
Call it $r_i(p)$.  Prefix compatibility in the two local lemmas implies
that computing with any larger terminal horizon gives the identical
rational entry.  To verify this assertion, fix the two finite runs and
induct on their common stage index: every lookup needed for a common
output row has the same input value by the induction hypothesis, so the
deterministic recursion returns the same value.  This finite argument establishes prefix consistency. For two output
precisions $p,q$, run the construction at terminal horizon $\max(p,q)$,
apply root regularity there, and identify its entries with the separately
defined outputs. Thus the outputs are regular names; no compatible-root
choice principle is needed.

\subsection{The uniform factorization theorem}

Set
\[
 G=1+\max_{i<d}A_i,\qquad
 W=1+\max_{i<d}D_i^+,\qquad
 c=\operatorname{bl}(1+dWG^d).
 \tag{F15}\label{eq:fta-final-budget}
\]

\begin{theorem}[Degree-uniform constructive factorization schema]
\label{thm:fta-uniform-factorization}
Fix one primitive recursive coefficient evaluator
$a(d,\mathbf p,i,s)$, with natural parameters $d,\mathbf p,i,s$.
There is one primitive recursive root evaluator
$r(d,M,\mathbf p,i,s)$, and the explicit primitive recursive bound
$c=c(d,M)$ above, such that iPRA proves the following.
If the input is monic of degree $d\ge1$, its coefficient names are
regular, and $|a(d,\mathbf p,i,s)|_1\le M$ for $i\le d$, then the
constructed $r_i$ are regular coefficient names and, for every $k$,
\[
 \left\|
   \sum_{j=0}^d a(d,\mathbf p,j,k+c)X^j
       -\prod_{i<d}(X-r_i(k+c))
 \right\|_1\le2^{-k}.
 \tag{F16}\label{eq:fta-uniform-factorization}
\]
In particular these names represent a full factorization, with
multiplicities.  All degrees are handled by the same numeric program,
not by a separately compiled evaluator for each degree.
\end{theorem}

\begin{proof}
Fix $k$ and put $t=k+c$.  Use the finite schedule
\eqref{eq:fta-backwards}.  The global input assumptions establish the
finite good-table condition at $h_0$ by bounded Boolean-fold induction.
Lemma~\ref{lem:fta-deflation-table} propagates the good-table conditions
through the $d$ stages.  Keep only the finite coefficients at the common
precision $t$.  Write $P_i$ for the degree-$d-i$ sampled polynomial,
$L_i=X-r_i(t)$, and $e_i=P_i-L_iP_{i+1}$.
Then $P_d=1$, $\|L_i\|_1\le G$, and
$\|e_i\|_1\le W2^{-t}$.  Repeated substitution gives the exact identity
\[
 P_0-\prod_{i<d}L_i
   =\sum_{i<d}\left(\prod_{j<i}L_j\right)e_i.
 \tag{F17}\label{eq:fta-telescope}
\]
Finite convolution bounds therefore give
\[
 \left\|P_0-\prod_{i<d}L_i\right\|_1
 \le W2^{-t}\sum_{i<d}G^i
 \le dWG^d2^{-t}\le2^{-k}.
\]
This proves \eqref{eq:fta-uniform-factorization}.

For the logical-strength claim, the stage invariant comprises the Boolean
good-table checks through the fixed $h_i$, the finitely many local
defect inequalities, and the finite coefficient identity
\eqref{eq:fta-telescope} through the current stage.  Every quantified
index is bounded by $d$, a listed $h_i$, or a coded array length.
Their Booleanization is quantifier-free.  Quantifier-free induction is
therefore enough.  Universal introduction over $k$ follows only after
this finite proof.  There is no induction on represented-complex equality
or on the assertion that a polynomial factors at all precisions.
\end{proof}

The dependence on the initial evaluator is exactly the dependence allowed
by the original FTA schema: once that PR definition is fixed, the table
generator is a composition with it.  For a finite rational-polynomial
code, constant coefficient names give a single first-order specialization
uniform in that code and its degree.  For an apart leading coefficient,
the usual certified reciprocal and a fixed precision shift reduce to the
monic formulation.  This does not decide which coefficient is leading
when the degree has not been supplied.

\subsection{Resource bounds and the remaining Lean work}

The query bound and verification bound have different roles.
The raw root at precision $t$ reads through $U(n,M,t)$; the normalized
root reads through $U(n,M,t+2)$.  A deflation through output horizon $h$
is verified from goodness through $V(n,M,h)$, which includes shifted
coefficient comparisons and the same-precision residual estimate.
Composing only the smaller query bound would not justify the theorem.

No non-primitive-recursive search is hidden in rational encoding sizes.
The fixed root table routine performs exactly $J(t)$ descent iterations;
each descent has bounded loops in its current finite codes, as spelled
out in the imported lemma.  The complete table map is therefore a fixed
PR function.  Iteration over $d$ takes table \emph{values} as state; it
never forms a growing PR program and asks a universal evaluator to run it.

If a numerical arithmetic-operation envelope is desired, first use the dyadic
normalization above.  A finite input table then has a primitive recursive
code bound $Z(n,M,h)$: with any fixed tuple encoding, bound the signed
dyadic numerators by $(M+1)2^{h+4}$ and the denominators by $2^{h+4}$,
then bound the $(n+1)(h+1)$ entries by a positional digit code.  Instrument
the one fixed table routine with an arithmetic-operation counter, using
the same compositions and recursions.  Its counter function $T$ is PR. This counts operations of the fixed table
transformer; the cost of generating the original coefficient table remains
the cost of its particular fixed evaluator. No bit-cost estimate is claimed.
The explicit bounded maximum
\[
 \max_{z\le Z(n,M,V(n,M,h))}
      T(n,M,h,z)
\]
is a PR envelope, with malformed tables totalized as above.  This maximum
is over finite input data for one fixed routine, not over arbitrary
programs.  An independent complexity claim stronger than primitive
recursiveness is not needed.

A Lean implementation should define the finite table data, $U,V$, the
cap recursion, rounding, and polynomial coefficient norm first.  The
critical new lemmas are the finite normalized-root contract,
the finite deflator contract, prefix compatibility, and the coefficient
telescoping bound.  Their precise statements are
Lemmas~\ref{lem:fta-root-table} and~\ref{lem:fta-deflation-table},
\eqref{eq:fta-telescope}, and
Theorem~\ref{thm:fta-uniform-factorization}.  Their arithmetic versions
must then construct object-language proofs using the Booleanized finite
invariants just described.  The repository's local descent derivation
remains an imported human proof until that object derivation is finished;
a Lean theorem that merely accepts the descent contract as a hypothesis
must be labelled conditional.
